\documentclass[12pt,a4paper]{article}
\usepackage[utf8]{inputenc}
\usepackage[spanish]{babel}
\usepackage{amsmath}
\usepackage{amssymb}
\usepackage{graphicx}
\usepackage{geometry}
\usepackage{listings}
\usepackage{xcolor}
\usepackage{booktabs}
\usepackage{hyperref}
\usepackage{float}
\usepackage{tikz}
\usepackage{fancyhdr}
\usepackage{setspace}

\geometry{margin=2.5cm}
\onehalfspacing

% Estilo para codigo Python
\definecolor{codegreen}{rgb}{0,0.6,0}
\definecolor{codegray}{rgb}{0.5,0.5,0.5}
\definecolor{codepurple}{rgb}{0.58,0,0.82}
\definecolor{backcolour}{rgb}{0.95,0.95,0.92}

\lstdefinestyle{pythonstyle}{
    backgroundcolor=\color{backcolour},
    commentstyle=\color{codegreen},
    keywordstyle=\color{blue},
    numberstyle=\tiny\color{codegray},
    stringstyle=\color{codepurple},
    basicstyle=\ttfamily\footnotesize,
    breakatwhitespace=false,
    breaklines=true,
    captionpos=b,
    keepspaces=true,
    numbers=left,
    numbersep=5pt,
    showspaces=false,
    showstringspaces=false,
    showtabs=false,
    tabsize=2,
    language=Python
}

\pagestyle{fancy}
\fancyhf{}
\rhead{Red Neuronal Artificial -- EMNIST Digits}
\lhead{Inteligencia Artificial}
\cfoot{\thepage}

\begin{document}

%----------------------------------------------------------------
% PORTADA
%----------------------------------------------------------------
\begin{titlepage}
    \centering
    \vspace*{1cm}
    {\Large \textbf{UNIVERSIDAD TÉCNICA DE AMBATO}}\\[0.5cm]
    {\large Facultad de Ingeniería en Sistemas, Electrónica e Industrial}\\[0.3cm]
    {\large Carrera de Software}\\[2cm]

    \rule{\linewidth}{0.5mm}\\[0.4cm]
    {\LARGE \textbf{Red Neuronal Artificial (ANN)\\para Clasificación de Dígitos\\con Dataset EMNIST}}\\[0.4cm]
    \rule{\linewidth}{0.5mm}\\[2cm]

    {\large \textbf{Materia:} Inteligencia Artificial}\\[0.3cm]
    {\large \textbf{Estudiante:} Angel Ayuquina}\\[0.3cm]
    {\large \textbf{Fecha:} Mayo 2026}\\[3cm]

    % [IMAGEN: Logo de la Universidad Técnica de Ambato]

    \vfill
\end{titlepage}

%----------------------------------------------------------------
% TABLA DE CONTENIDOS
%----------------------------------------------------------------
\tableofcontents
\newpage

%----------------------------------------------------------------
% 1. INTRODUCCIÓN
%----------------------------------------------------------------
\section{Introducción}

Las redes neuronales artificiales (ANN, por sus siglas en inglés \textit{Artificial Neural Networks}) constituyen uno de los pilares fundamentales del aprendizaje automático moderno. Inspiradas en el funcionamiento del cerebro humano, estas estructuras computacionales son capaces de aprender patrones complejos a partir de grandes volúmenes de datos.

En el presente trabajo se implementa una red neuronal artificial multicapa utilizando el framework \textbf{TensorFlow/Keras} con el lenguaje de programación Python, siguiendo la metodología explicada en clase. El objetivo principal es entrenar la red para que \textbf{reconozca dígitos escritos a mano} (del 0 al 9) utilizando el dataset \textbf{EMNIST Digits}.

\subsection{Objetivos}
\begin{itemize}
    \item Implementar una red neuronal artificial usando \texttt{keras.Sequential()}.
    \item Aplicar los conceptos de función de activación, función de pérdida y optimizador Adam.
    \item Entrenar el modelo con datos de entrenamiento durante 30 épocas.
    \item Evaluar el modelo con datos de testeo y calcular el error de testeo.
    \item Demostrar que el error de testeo cumple con el umbral establecido ($\leq 0.10$).
\end{itemize}

\subsection{Librerías utilizadas}
Por indicación del docente, únicamente se utilizaron las siguientes librerías:
\begin{itemize}
    \item \textbf{NumPy}: manejo de arrays y cálculo del error de testeo.
    \item \textbf{TensorFlow}: framework de aprendizaje profundo.
    \item \textbf{Keras}: API de alto nivel para construir y entrenar la red neuronal.
\end{itemize}

%----------------------------------------------------------------
% 2. MARCO TEÓRICO
%----------------------------------------------------------------
\section{Marco Teórico}

\subsection{Red Neuronal Artificial}

Una red neuronal artificial es un modelo computacional compuesto por capas de nodos (neuronas) interconectadas. Cada neurona recibe entradas, las pondera con pesos, aplica una función de activación y produce una salida. Matemáticamente, la operación de una neurona se expresa como:

\begin{equation}
    y = f\left(\sum_{i=1}^{n} w_i x_i + b\right)
\end{equation}

donde:
\begin{itemize}
    \item $x_i$ son las entradas
    \item $w_i$ son los pesos sinápticos
    \item $b$ es el sesgo (\textit{bias})
    \item $f(\cdot)$ es la función de activación
    \item $y$ es la salida de la neurona
\end{itemize}

\subsection{Modelo Secuencial (keras.Sequential)}

El modelo \texttt{keras.Sequential} permite construir redes neuronales capa por capa de forma lineal. Cada capa recibe como entrada la salida de la capa anterior. Solo la primera capa requiere el parámetro \texttt{input\_shape}, que define el número de características (columnas) del dataset de entrada.

\subsection{Funciones de Activación}

Las funciones de activación introducen no linealidad en la red, permitiendo aprender patrones complejos. Las utilizadas en este proyecto son:

\subsubsection{ReLU (Rectified Linear Unit)}
\begin{equation}
    f(x) = \max(0, x)
\end{equation}
Es la función más utilizada en capas ocultas. Resuelve el problema del \textit{vanishing gradient} (desvanecimiento de gradientes) ya que su derivada es constante para valores positivos. Como explicó el docente: \textit{"los pesos entre 0 y 1 son tan pequeñitos que ya no aprenden nada, entonces ahí se utiliza un arreglo"} -- ese arreglo es precisamente ReLU.

% [IMAGEN: Gráfica de la función ReLU mostrando f(x)=0 para x<0 y f(x)=x para x>=0]

\subsubsection{Tanh (Tangente Hiperbólica)}
\begin{equation}
    f(x) = \frac{e^x - e^{-x}}{e^x + e^{-x}}
\end{equation}
Produce salidas en el rango $[-1, 1]$. El docente indicó que \textit{"cada capa puede tener una distinta"} función de activación, lo que justifica su uso en la tercera capa oculta.

\subsubsection{Sigmoid}
\begin{equation}
    f(x) = \frac{1}{1 + e^{-x}}
\end{equation}
Produce salidas en el rango $(0, 1)$. Se utiliza en la cuarta capa oculta, demostrando el concepto de activaciones distintas por capa.

\subsubsection{Softmax}
\begin{equation}
    f(x_i) = \frac{e^{x_i}}{\sum_{j=1}^{K} e^{x_j}}
\end{equation}
Convierte las salidas en probabilidades que suman 1. Se utiliza \textbf{exclusivamente en la capa de salida} para clasificación multiclase (10 dígitos). La neurona con mayor probabilidad determina el dígito predicho.

\subsection{Función de Pérdida: Entropía Cruzada}

Para problemas de clasificación multiclase se utiliza la \textbf{entropía cruzada categórica dispersa} (\textit{sparse categorical crossentropy}):

\begin{equation}
    \mathcal{L} = -\sum_{i=1}^{N} \log\left(\hat{p}_{i,y_i}\right)
\end{equation}

donde $\hat{p}_{i,y_i}$ es la probabilidad predicha para la clase verdadera $y_i$ de la muestra $i$. El objetivo del entrenamiento es \textbf{minimizar} esta función de pérdida.

\subsection{Optimizador Adam}

Adam (\textit{Adaptive Moment Estimation}) es un algoritmo de optimización basado en el descenso del gradiente estocástico. Como explicó el docente con la analogía de la montaña:

\begin{quote}
\textit{"Me dejan en cualquier parte de la montaña con los ojos vendados y me dicen: desciende a lo más bajo. Al inicio puedo dar saltos grandes; cuando está llegando al mínimo global, da saltos pequeños."}
\end{quote}

Adam actualiza los pesos mediante:

\begin{equation}
    w_{t+1} = w_t - \frac{\alpha}{\sqrt{\hat{v}_t} + \epsilon} \cdot \hat{m}_t
\end{equation}

donde $\hat{m}_t$ es el primer momento (media del gradiente), $\hat{v}_t$ es el segundo momento (varianza del gradiente) y $\alpha$ es la tasa de aprendizaje.

\subsection{Retropropagación del Error}

Durante el entrenamiento, el error se propaga hacia atrás a través de la red para actualizar los pesos:

\begin{equation}
    \frac{\partial \mathcal{L}}{\partial w_{ij}} = \frac{\partial \mathcal{L}}{\partial a_j} \cdot \frac{\partial a_j}{\partial z_j} \cdot \frac{\partial z_j}{\partial w_{ij}}
\end{equation}

Este proceso se repite en cada época de entrenamiento.

%----------------------------------------------------------------
% 3. DATASET EMNIST DIGITS
%----------------------------------------------------------------
\section{Dataset: EMNIST Digits}

\subsection{Descripción}

El dataset \textbf{EMNIST Digits} (\textit{Extended Modified National Institute of Standards and Technology}) es una extensión del famoso dataset MNIST, diseñado para el reconocimiento de dígitos escritos a mano.

% [IMAGEN: Ejemplos de imágenes del dataset EMNIST Digits mostrando dígitos del 0 al 9]

\subsection{Características del dataset}

\begin{table}[H]
    \centering
    \begin{tabular}{@{}ll@{}}
        \toprule
        \textbf{Característica} & \textbf{Valor} \\
        \midrule
        Tipo de datos & Imágenes en escala de grises \\
        Tamaño de cada imagen & 28 × 28 píxeles \\
        Número de características (n\_features) & 784 píxeles por imagen \\
        Clases & 10 (dígitos del 0 al 9) \\
        Muestras de entrenamiento & 240,000 \\
        Muestras de testeo & 40,000 \\
        Formato & CSV (columna 0: etiqueta, columnas 1-784: píxeles) \\
        \bottomrule
    \end{tabular}
    \caption{Características del dataset EMNIST Digits}
\end{table}

\subsection{Preprocesamiento}

\subsubsection{Separación de etiquetas y características}
Cada fila del CSV contiene 785 valores: la primera columna es la etiqueta (dígito real 0-9) y las 784 columnas restantes son los valores de píxeles.

\begin{lstlisting}[style=pythonstyle, caption=Carga y separación del dataset]
datos_train = np.loadtxt(RUTA_TRAIN, delimiter=',', dtype='float32')

y_train = datos_train[:, 0].astype(int)   # etiquetas
X_train = datos_train[:, 1:].astype('float32')  # pixeles

n_features = X_train.shape[1]  # = 784
\end{lstlisting}

\subsubsection{Normalización}
Los valores de píxeles originalmente están en el rango $[0, 255]$. Se normalizan al rango $[0.0, 1.0]$ dividiendo entre 255:

\begin{equation}
    X_{norm} = \frac{X}{255}
\end{equation}

\begin{lstlisting}[style=pythonstyle, caption=Normalización de píxeles]
X_train = X_train / 255.0
X_test  = X_test  / 255.0
\end{lstlisting}

Esta normalización mejora la convergencia del descenso del gradiente al mantener los valores de entrada en una escala pequeña y uniforme.

%----------------------------------------------------------------
% 4. ARQUITECTURA DE LA RED NEURONAL
%----------------------------------------------------------------
\section{Arquitectura de la Red Neuronal}

\subsection{Diseño de la arquitectura}

La red neuronal se construye usando \texttt{keras.Sequential()}, instanciado en una variable llamada \texttt{model}. La arquitectura diseñada demuestra el concepto explicado por el docente: \textit{"cada capa puede tener una función de activación distinta"}.

\begin{table}[H]
    \centering
    \begin{tabular}{@{}ccccc@{}}
        \toprule
        \textbf{Capa} & \textbf{Tipo} & \textbf{Neuronas} & \textbf{Activación} & \textbf{Parámetros} \\
        \midrule
        Entrada & -- & 784 & -- & -- \\
        Oculta 1 & Dense & 128 & ReLU & $784 \times 128 + 128 = 100,480$ \\
        Oculta 2 & Dense & 64 & ReLU & $128 \times 64 + 64 = 8,256$ \\
        Oculta 3 & Dense & 32 & Tanh & $64 \times 32 + 32 = 2,080$ \\
        Oculta 4 & Dense & 16 & Sigmoid & $32 \times 16 + 16 = 528$ \\
        Salida & Dense & 10 & Softmax & $16 \times 10 + 10 = 170$ \\
        \midrule
        \multicolumn{4}{l}{\textbf{Total de parámetros entrenables}} & \textbf{111,514} \\
        \bottomrule
    \end{tabular}
    \caption{Arquitectura de la red neuronal}
\end{table}

\subsection{Cálculo de parámetros}

El número de parámetros de cada capa densa se calcula como:
\begin{equation}
    \text{Parámetros} = (\text{entradas} \times \text{neuronas}) + \text{neuronas (bias)}
\end{equation}

Por ejemplo, para la primera capa oculta:
\begin{equation}
    \text{Parámetros}_1 = (784 \times 128) + 128 = 100{,}352 + 128 = 100{,}480
\end{equation}

% [IMAGEN: Diagrama de la arquitectura de la red neuronal mostrando las 5 capas con sus neuronas y activaciones, generado por keras.utils.plot_model o TensorBoard GRAPHS]

\subsection{Código de la arquitectura}

\begin{lstlisting}[style=pythonstyle, caption=Construcción del modelo con keras.Sequential]
model = keras.Sequential()

# Solo la primera capa lleva input_shape (n_features=784)
model.add(layers.Dense(128, activation='relu', input_shape=(n_features,)))

# Cada capa puede tener distinta activacion (como explico el profesor)
model.add(layers.Dense(64,  activation='relu'))
model.add(layers.Dense(32,  activation='tanh'))
model.add(layers.Dense(16,  activation='sigmoid'))

# Salida: 10 neuronas (una por digito 0-9), softmax para probabilidades
model.add(layers.Dense(10, activation='softmax'))
\end{lstlisting}

%----------------------------------------------------------------
% 5. COMPILACIÓN Y ENTRENAMIENTO
%----------------------------------------------------------------
\section{Compilación y Entrenamiento del Modelo}

\subsection{Compilación}

Una vez definida la arquitectura, el modelo se compila con \texttt{model.compile()}:

\begin{lstlisting}[style=pythonstyle, caption=Compilación del modelo]
model.compile(
    loss='sparse_categorical_crossentropy',
    optimizer='adam',
    metrics=['accuracy']
)
\end{lstlisting}

Los tres parámetros clave son:
\begin{itemize}
    \item \textbf{loss}: función de pérdida -- entropía cruzada categórica.
    \item \textbf{optimizer}: Adam, que gestiona la tasa de aprendizaje adaptativamente.
    \item \textbf{metrics}: exactitud (\textit{accuracy}) para monitorear el entrenamiento.
\end{itemize}

El \textbf{error de testeo} se obtiene como:
\begin{equation}
    \text{error} = 1 - \text{exactitud}
\end{equation}

\subsection{Entrenamiento}

El entrenamiento se ejecuta con \texttt{model.fit()} durante \textbf{30 épocas}, como indicó el docente:

\begin{lstlisting}[style=pythonstyle, caption=Entrenamiento del modelo]
historial = model.fit(
    X_train,
    y_train,
    epochs=30,
    batch_size=64,
    verbose=1
)
\end{lstlisting}

\subsection{Evolución del entrenamiento}

La siguiente tabla muestra la exactitud obtenida en cada época durante el entrenamiento:

\begin{table}[H]
    \centering
    \begin{tabular}{@{}cccc@{}}
        \toprule
        \textbf{Época} & \textbf{Accuracy} & \textbf{Época} & \textbf{Accuracy} \\
        \midrule
        1  & 89.68\% & 16 & 98.01\% \\
        2  & 95.83\% & 17 & 98.06\% \\
        3  & 96.57\% & 18 & 98.08\% \\
        4  & 96.98\% & 19 & 98.12\% \\
        5  & 97.20\% & 20 & 98.11\% \\
        6  & 97.37\% & 21 & 98.17\% \\
        7  & 97.52\% & 22 & 98.18\% \\
        8  & 97.63\% & 23 & 98.19\% \\
        9  & 97.71\% & 24 & 98.22\% \\
        10 & 97.76\% & 25 & 98.24\% \\
        11 & 97.80\% & 26 & 98.24\% \\
        12 & 97.85\% & 27 & 98.25\% \\
        13 & 97.92\% & 28 & 98.27\% \\
        14 & 97.96\% & 29 & 98.29\% \\
        15 & 98.00\% & 30 & 98.30\% \\
        \bottomrule
    \end{tabular}
    \caption{Exactitud del modelo por época durante el entrenamiento}
\end{table}

% [IMAGEN: Gráfica de TensorBoard mostrando la curva epoch_accuracy subiendo de 89.68% a 98.30% a lo largo de las 30 épocas. Captura de pantalla de la pestaña SCALARS en http://localhost:6006]

% [IMAGEN: Gráfica de TensorBoard mostrando la curva epoch_loss bajando a lo largo de las 30 épocas, demostrando la convergencia del modelo]

\textbf{Observación}: En la época 1 la exactitud comienza en 89.68\%, evidenciando que el modelo inicia con pesos aleatorios. A partir de la época 5 supera el 97\% y converge progresivamente hacia 98.30\% en la época 30.

%----------------------------------------------------------------
% 6. EVALUACIÓN CON DATOS DE TESTEO
%----------------------------------------------------------------
\section{Evaluación con Datos de Testeo}

\subsection{Proceso de predicción}

Como explicó el docente, las etiquetas de testeo (\texttt{y\_test}) se guardan aparte desde el inicio y \textbf{no se utilizan durante el entrenamiento}. Tras el entrenamiento, se ejecuta \texttt{model.predict()} exclusivamente con \texttt{X\_test}:

\begin{lstlisting}[style=pythonstyle, caption=Predicción y evaluación con datos de testeo]
# Guardar y_test aparte (etiquetas verdaderas)
y_verdaderas = y_test

# model.predict() SOLO con X_test
y_pred_probabilidades = model.predict(X_test)

# La neurona con mayor probabilidad = clase predicha
y_predichas = np.argmax(y_pred_probabilidades, axis=1)

# Calcular error de testeo
correctas        = np.sum(y_predichas == y_verdaderas)
total            = len(y_verdaderas)
exactitud_testeo = correctas / total
error_testeo     = 1 - exactitud_testeo
\end{lstlisting}

\subsection{Resultados obtenidos}

\begin{table}[H]
    \centering
    \begin{tabular}{@{}ll@{}}
        \toprule
        \textbf{Métrica} & \textbf{Valor} \\
        \midrule
        Total de muestras de testeo & 40,000 \\
        Predicciones correctas & 39,049 \\
        Exactitud de testeo & 97.62\% \\
        Error de testeo & 0.0238 \\
        Umbral máximo de error (docente) & 0.10 \\
        \midrule
        \textbf{Veredicto} & \textbf{Arquitectura ADECUADA} \\
        \bottomrule
    \end{tabular}
    \caption{Resultados de la evaluación con datos de testeo}
\end{table}

\subsection{Verificación del umbral}

El docente estableció que el error de testeo \textbf{no debe ser mayor a 0.10}. El modelo obtuvo:

\begin{equation}
    \text{error\_testeo} = 1 - 0.9762 = 0.0238 \leq 0.10 \quad \checkmark
\end{equation}

Esto confirma que la arquitectura diseñada es \textbf{adecuada} para la tarea de clasificación de dígitos.

\subsection{Comparación: entrenamiento vs testeo}

\begin{table}[H]
    \centering
    \begin{tabular}{@{}lcc@{}}
        \toprule
        \textbf{Conjunto} & \textbf{Exactitud} & \textbf{Error} \\
        \midrule
        Entrenamiento (época 30) & 98.30\% & 0.0170 \\
        Testeo & 97.62\% & 0.0238 \\
        \midrule
        Diferencia & 0.68\% & 0.0068 \\
        \bottomrule
    \end{tabular}
    \caption{Comparación de exactitud entre entrenamiento y testeo}
\end{table}

La diferencia mínima entre entrenamiento (98.30\%) y testeo (97.62\%) indica que el modelo \textbf{generaliza bien} y no presenta sobreajuste significativo (\textit{overfitting}).

%----------------------------------------------------------------
% 7. ESTRUCTURA DEL PROYECTO
%----------------------------------------------------------------
\section{Estructura del Proyecto}

El proyecto se organizó en carpetas separadas, cada una con una responsabilidad específica:

\begin{lstlisting}[language=bash, caption=Estructura de carpetas del proyecto]
red_neuronal_emnist/
|
|-- main.py                    <- Punto de entrada principal
|
|-- datos/
|   |-- cargar_datos.py        <- Carga CSV y normaliza pixeles
|
|-- modelo/
|   |-- arquitectura.py        <- keras.Sequential + model.add()
|   |-- compilar.py            <- model.compile() + model.summary()
|   |-- entrenar.py            <- model.fit() con 30 epocas
|
|-- evaluacion/
|   |-- testeo.py              <- model.predict() + error de testeo
|
|-- visualizacion/
|   |-- tensorboard_viz.py     <- TensorBoard para graficas
|   |-- ver_red.py             <- Diagrama de la arquitectura
\end{lstlisting}

% [IMAGEN: Captura de pantalla de VS Code mostrando la estructura de carpetas del proyecto en el explorador de archivos]

%----------------------------------------------------------------
% 8. CONCLUSIONES
%----------------------------------------------------------------
\section{Conclusiones}

\begin{enumerate}
    \item Se implementó exitosamente una red neuronal artificial multicapa usando \texttt{keras.Sequential()} siguiendo exactamente la metodología explicada por el docente.

    \item La arquitectura con capas de activaciones distintas (ReLU, Tanh, Sigmoid, Softmax) demostró el concepto impartido en clase: \textit{"cada capa puede tener una función de activación distinta"}.

    \item El modelo alcanzó una \textbf{exactitud de testeo del 97.62\%} con un \textbf{error de 0.0238}, ampliamente por debajo del umbral máximo de 0.10 establecido por el docente.

    \item El optimizador Adam gestionó eficientemente la tasa de aprendizaje durante el descenso del gradiente, permitiendo que el modelo convergiera en 30 épocas.

    \item La normalización de píxeles al rango $[0, 1]$ fue fundamental para acelerar la convergencia del entrenamiento.

    \item La diferencia mínima entre la exactitud de entrenamiento (98.30\%) y testeo (97.62\%) indica que el modelo generaliza correctamente sin sobreajuste significativo.
\end{enumerate}

%----------------------------------------------------------------
% 9. REFERENCIAS
%----------------------------------------------------------------
\section{Referencias}

\begin{itemize}
    \item Cohen, G., Afshar, S., Tapson, J., \& van Schaik, A. (2017). EMNIST: an extension of MNIST to handwritten letters. \textit{arXiv preprint arXiv:1702.05373}.

    \item Chollet, F. (2021). \textit{Deep Learning with Python} (2nd ed.). Manning Publications.

    \item TensorFlow Development Team. (2024). \textit{TensorFlow 2.x Documentation}. Recuperado de \url{https://www.tensorflow.org/}

    \item Keras Development Team. (2024). \textit{Keras Documentation}. Recuperado de \url{https://keras.io/}

    \item Goodfellow, I., Bengio, Y., \& Courville, A. (2016). \textit{Deep Learning}. MIT Press.
\end{itemize}

\end{document}
